\scanpage{25}\mspage{24}
\begingroup\csname tagsleft@true\endcsname

Je reviens sur la description \unsure{d'un|d'une} atelier spacieux ou polyédral, satisfaisant L1, L2, L3 plus \deleted{At \unsure{L(spac)|L1(spac)} ou} At CL(spac) (ou, au choix At CL(pol)), en termes de $(M,\leq,\ll,\mathrel{>\mkern-10mu<})^{*}$.
% slanted
\margin{* ce qui détermine \unsure{tout} au moins les \deleted{\unsure{familles}} figures de \unsure{réf.} \deleted{\unsure{ou|et}} \unsure{et} les sous-figures de telles\ldots}

Le point ici, c'est que le critère At CL(spac) réduit, pour connaître \({\mathrel{>\mkern-10mu<}}_M\), à la connaissance de $\deleted{\mathrel{>\mkern-10mu<}}{\Vert}_{\mathcal L}$, i.e. à celle d'une structure de graphe sur $\mathcal L$. (NB ds le cas d'un atelier ensembliste, je rappelle que $x\mathrel{\Vert}y\Longleftrightarrow x\neq y$)

Donc, la question est, partant de \deleted{$(M,\leq,\ll)$}, \deleted{\unsure{d'où|de}, \unsure{et} $\mathcal L$.}
\[
\unsure{,}\quad(M,\leq,\ll)\tag*{(1)}
\]
d'où
\[
\mathcal L=\{x\in M\mid x\text{ minimal pour }\ll\},\tag*{(2)}
\]
et d'une relation symétrique antiréflexive
\[
x\mathrel{\Vert}y\qquad(x,y\in\mathcal L)
\]
(\unsure{impliquant} $x\neq y$) sur $\mathcal L$, quelles conditions imposer pour que ça corresponde bien à un atelier\textsuperscript{\#} spacieux, satisfaisant \deleted{\unsure{At}} (en plus At \deleted{\unsure{L}}CL(spac)) les conditions L1, L2, L3.
% slanted
\margin{\# Il est entendu qu'on prendra pour $\unsure{F|\mathscr F}$, soit l'ens. de \uline{toutes} les \deleted{familles}, parties fermées de $M_{\leq}$ formées d'\unsure{él\textsuperscript{ts}|él\textsuperscript{t}} deux à deux compatibles, soit \unsure{qu'on précise celles|qu'on précise les} qui sont des s.f. (en supposant que tt partie fermée d'un $\widetilde X$ en est \unsure{une s.f.|de s.f.}).}

Il est entendu qu'on va \uline{définir} \({\mathrel{>\mkern-10mu<}}_M\) par
\[
\begin{aligned}
X\unsure{,|{}_{1}}\mathrel{>\mkern-10mu<}Y\quad&\Longleftrightarrow\quad
\forall\,\begin{aligned}[t]&x\in\operatorname{omb}(X)\setminus\operatorname{omb}(L)\\
&y\in\operatorname{omb}(Y)-\operatorname{omb}(L)\end{aligned},\\
&\hspace{5em}\text{on a}\quad x\mathrel{\Vert}y
\end{aligned}\tag*{(3)}
\]
\margin{\diagram{\resizebox{\linewidth}{!}{\begin{tikzcd}[ampersand replacement=\&,row sep=0.25em,cells={nodes={anchor=west}},
execute at end picture={
\draw ([xshift=-3pt]where.south west) .. controls +(right:.4em) and +(up:.8em) .. ([xshift=-1pt]bound.west) .. controls +(down:.8em) and +(left:.6em) .. (bound.south east);
}]
|[alias=where]| \text{où}\quad L=X\cap Y\\
\phantom{\text{où}\quad L}=\underset{Z\in\widetilde X\cap\widetilde Y}{\operatorname{Sup}}Z\\
\unsure{\text{donc}|\text{d'où}|\text{dans ce}}\\
|[alias=bound]| \operatorname{omb}(L)\overset{\text{déf}}{=}\underset{Z\in\widetilde X\cap\widetilde Y}{\bigcup}\operatorname{omb}(Z\unsure{)}
\end{tikzcd}}}}
Pour commencer, il faut \deleted{\unsure{vérifier}} \inserted{écrire} la condition préliminaire
\[
X\mathrel{>\mkern-10mu<}Y,\ X'\leq X,\ Y'\leq Y\ \Longrightarrow\ X'\mathrel{>\mkern-10mu<}Y'\tag*{(4)}
\]
\deleted{et on a besoin ce faisant \unsure{de|en} sous \unsure{formes|forme} plus}
\endgroup
